\section*{摘要}
\addcontentsline{toc}{section}{摘要}
本项目面向端侧实时绘画场景，在易灵思Ti60F225I3平台构建摄像头采集、神经网络风格化与HDMI显示系统。SC431HAI的四lane RAW10经协议检查、Bayer插值和尺寸变换形成1280×720 RGB565图像，经DDR交给六通道轻量卷积加速器。网络由输入投影、四组膨胀DW/PW残差块及有符号RGB输出头组成，各层保持全分辨率。绘画Gram统计、教师引导、局部纹理约束与整数QAT共同用于训练十组独立风格参数。

RISC-V负责摄像头初始化、216项参数装载、作业与中断控制。三源、三输出槽通过完整帧发布和安全消隐提交协调并发，支持原图、风格图及同源对比；按键控制视图、风格和UI，当前板载手势采用个人模板及四秒确认。最终CAD资源为230块RAM10和142个DSP，最小setup/hold余量为+0.108/+0.026ns。完整720p帧的CNN/DMA仿真耗时25.585ms，HDMI输出刷新率为60Hz。

\noindent\textbf{关键词：}FPGA；神经风格迁移；深度可分离卷积；定点量化；RISC-V；视频交互
